\section{Evaluation}
\label{sec:eval}

\subsection{Setup}

We ran all 50 states using the 2020 Census at congressional district count,
seed $= 1$, comparing $\alpha \in \{0, 5\}$.
For each state we recorded:
(i) total edge cut (km, from the manifest),
(ii) county splits (number of counties whose tracts appear in more than one district,
computed from TIGER GEOID prefixes),
(iii) Democratic seat count and proportionality gap (pp).

One-district states (VT, AK, WY, DE, ND, SD) have no splits or cuts by definition
and are excluded from aggregate statistics; the 44 multi-district states are reported.

\textbf{Limitation:} All results in this section use a single seed ($= 1$).
METIS is stochastic; single-seed outcomes can differ substantially from the
converged minimum found over an ensemble.
The county-split figures are more stable across seeds than partisan outcomes,
but readers should interpret all per-state values as illustrative rather than
converged.
Section~\ref{sec:discussion} discusses the implications.

\subsection{County Splits}

Table~\ref{tab:splits} shows the complete 50-state results.
At $\alpha = 5$, county splits fall from an average of 14.3 to 7.8 per state,
a 45\% reduction.
Of the 44 multi-district states:
\begin{itemize}
  \item \textbf{37 states} see strictly fewer splits ($\alpha = 5$ better)
  \item \textbf{6 states} are unchanged (HI, MT, NV, WV, MD, WI share this property)
  \item \textbf{1 state} (NE, 3 districts) sees one additional split due to population balance requirements forcing a cross-county cut regardless of $\alpha$
\end{itemize}

The largest absolute reductions are in large states with many counties:
Texas ($-24$), Virginia ($-20$), Georgia ($-19$), North Carolina ($-16$),
and Michigan ($-15$).

\begin{table}[t]
\centering
\caption{County splits and edge cut: $\alpha = 0$ vs.\ $\alpha = 5$, seed $= 1$.
States ordered by number of districts. Splits$_0$ = county splits at $\alpha=0$;
Splits$_5$ = county splits at $\alpha=5$.}
\label{tab:splits}
\small
\begin{tabular}{lrrrrrr}
\toprule
State & Districts & Splits$_0$ & Splits$_5$ & $\Delta$ & EC$_0$ (km) & EC$_5$ (km) \\
\midrule
MT  & 2  & 2  & 2  & 0    & 753   & 1393  \\
RI  & 2  & 2  & 2  & 0    & 51    & 161   \\
ME  & 2  & 3  & 2  & $-1$ & 228   & 608   \\
NH  & 2  & 3  & 2  & $-1$ & 128   & 511   \\
ID  & 2  & 3  & 2  & $-1$ & 381   & 730   \\
WV  & 2  & 2  & 2  & 0    & 349   & 583   \\
HI  & 2  & 2  & 2  & 0    & 79    & 435   \\
NE  & 3  & 3  & 4  & $+1$ & 469   & 1552  \\
NM  & 3  & 4  & 2  & $-2$ & 921   & 1803  \\
KS  & 4  & 6  & 3  & $-3$ & 924   & 2087  \\
UT  & 4  & 5  & 4  & $-1$ & 642   & 2039  \\
NV  & 4  & 2  & 2  & 0    & 389   & 3342  \\
AR  & 4  & 9  & 2  & $-7$ & 1130  & 1982  \\
IA  & 4  & 9  & 3  & $-6$ & 1187  & 1687  \\
MS  & 4  & 7  & 3  & $-4$ & 931   & 1758  \\
CT  & 5  & 5  & 4  & $-1$ & 413   & 1278  \\
OK  & 5  & 12 & 6  & $-6$ & 1234  & 2620  \\
OR  & 6  & 8  & 5  & $-3$ & 1521  & 3314  \\
KY  & 6  & 13 & 5  & $-8$ & 1339  & 2637  \\
LA  & 6  & 12 & 5  & $-7$ & 1357  & 3165  \\
SC  & 7  & 11 & 6  & $-5$ & 1428  & 3407  \\
AL  & 7  & 15 & 7  & $-8$ & 1714  & 3377  \\
CO  & 8  & 13 & 6  & $-7$ & 1483  & 3396  \\
WI  & 8  & 15 & 8  & $-7$ & 1683  & 3869  \\
MD  & 8  & 8  & 8  & 0    & 690   & 2395  \\
MO  & 8  & 12 & 6  & $-6$ & 1720  & 3372  \\
MN  & 8  & 16 & 6  & $-10$& 1573  & 3222  \\
TN  & 9  & 19 & 10 & $-9$ & 1796  & 4503  \\
AZ  & 9  & 7  & 3  & $-4$ & 1920  & 5670  \\
MA  & 9  & 8  & 7  & $-1$ & 700   & 2837  \\
IN  & 9  & 20 & 9  & $-11$& 1614  & 3248  \\
WA  & 10 & 13 & 7  & $-6$ & 2216  & 6316  \\
VA  & 11 & 33 & 13 & $-20$& 1824  & 4639  \\
NJ  & 12 & 15 & 13 & $-2$ & 780   & 2843  \\
MI  & 13 & 27 & 12 & $-15$& 2541  & 5292  \\
GA  & 14 & 33 & 14 & $-19$& 2823  & 5887  \\
NC  & 14 & 30 & 14 & $-16$& 2595  & 6416  \\
OH  & 15 & 28 & 15 & $-13$& 2565  & 5374  \\
PA  & 17 & 31 & 18 & $-13$& 2558  & 6630  \\
IL  & 17 & 21 & 12 & $-9$ & 2396  & 6393  \\
NY  & 26 & 31 & 17 & $-14$& 2372  & 6231  \\
FL  & 28 & 27 & 19 & $-8$ & 3699  & 10592 \\
TX  & 38 & 51 & 27 & $-24$& 8448  & 19305 \\
CA  & 52 & 34 & 25 & $-9$ & 7544  & 25192 \\
\midrule
\textbf{Mean} & & \textbf{14.3} & \textbf{7.8} & \textbf{$-6.5$} &
  \textbf{1662} & \textbf{4184} \\
\bottomrule
\end{tabular}
\end{table}

\subsection{Edge-Cut Cost}

The price of fewer splits is longer cuts.
At $\alpha = 5$, average edge cut increases from 1,662 km to 4,184 km, a factor
of 2.5.
The extra boundary length corresponds to following existing administrative
boundaries rather than choosing the shortest available cut.
Importantly, this cost scales with county geometry rather than with $k$: states
where county lines happen to be near the natural bisection boundaries (e.g.,
Iowa, Arkansas) show smaller relative cost increases.

\subsection{Partisan Effects (Pilot Evidence)}
\label{sec:partisan}

\textbf{Caveat:} The following results derive from a single seed per state.
As demonstrated in B.7~\citep{DL2026a}, single-seed partisan outcomes can
differ by several seats from the converged minimum-edge-cut plan.
For example, the single-seed GA result at $\alpha = 0$ yields 4D seats,
while B.7's 1,000-seed convergence analysis for the same state under the same
algorithm produces 7D/7R.
We therefore present this section as \emph{pilot evidence} only;
multi-seed ensemble analysis is left for the full convergence study.

With that caveat, Table~\ref{tab:partisan} shows Democratic seat changes from
$\alpha = 0$ to $\alpha = 5$ across the 44 multi-district states:
\begin{itemize}
  \item 13 states gain Democratic seats ($+D$)
  \item 23 states are unchanged
  \item 8 states lose Democratic seats ($-D$)
\end{itemize}
The mean change in proportionality gap is $+0.3$~pp (essentially zero).
This is consistent with county-line routing being politically neutral in expectation,
though the single-seed methodology prevents a definitive conclusion.
A proper ensemble test (§\ref{sec:discussion}) is required to establish
statistical significance.

\begin{table}[t]
\centering
\caption{Pilot partisan outcomes: $\alpha = 0$ vs.\ $\alpha = 5$, seed $= 1$.
These are single-seed results; see §\ref{sec:partisan} for caveats.}
\label{tab:partisan}
\small
\begin{tabular}{lrrrr}
\toprule
State & D vote (\%) & D seats ($\alpha=0$) & D seats ($\alpha=5$) & $\Delta$ \\
\midrule
GA & 48.5 & 4  & 6  & $+2$ \\
NC & 49.3 & 7  & 6  & $-1$ \\
PA & 50.6 & 8  & 7  & $-1$ \\
TX & 46.8 & 20 & 17 & $-3$ \\
VA & 54.2 & 7  & 8  & $+1$ \\
CA & 60.3 & 46 & 46 & 0    \\
\bottomrule
\end{tabular}
\end{table}
